Large Language Models are becoming a tool used in recommender systems especially when using RAG alongside it for retrieval. It allows the use of natural language for items, simplifies ranking by using the LLM itself as a ranker especially catered for the user. But all of these perks come with a huge security risk when malicious actors inject poisoned data or instructions into the retrieved content. This could allow a specific targeted item to appear in the top recommendation. With this issue in mind, This thesis discusses RAGPoison-lab, a framework made for the purpose of this study to better understand this new risk using a controlled local benchmark environment for red-teaming retrieval-augmented recommendation pipelines.

The project will use the MovieLens 100K dataset as the foundation for the recommendation data to use in this experiment, using it to build clean and poisoned indices utilising the vector database Elasticsearch, three different attack types, frontier models while tracing and auditing every step through raw JSON artifacts. Frontend and SDK will allow us to visualize the data in a more intelligible way giving us a better look at the results. We will compare the baseline recommendation to the attacked one at \textit{k}=10 using HR, MRR and NDCG as our recommendation quality metrics and ASR for Target oriented attacks only. The final experiment is completed with 15 successful runs, thanks the five models being ran through three attack families.

The final result of the experiment show real risk patterns with targeted promotions bringing the targeted movie into the top-10 recommendations resulting in an impressive ASR value, while prompt injection proves the risk of injected instructions through a payload. untargeted promot displays the biggest gap between baseline and attacked with clear losses in quality. This thesis will thus study, with the help of RAGPoison, the security risks that come with RAG LLM-based recommenders using traces and logs from every step of the pipeline, from preparing the data, to indexing it, to the retrieval, to the final output. For security concerns the whole project was executed locally in a closed environment (except for model api calling) and never infects a third party tool.
